\documentclass[11pt]{article}
\usepackage{amsmath,amsfonts}
\usepackage{graphicx}
%\usepackage{xcolor}
%\definecolor{gray}{rgb}{.75,.75,.75}
%\usepackage[landscape]{geometry}
%\usepackage{hyperref}
%\usepackage[bookmarks=true, pdfborder={0 0 .5}, urlbordercolor=gray]{hyperref}

\renewcommand{\thefootnote}{\fnsymbol{footnote}}


\addtolength{\topmargin}{-.75in}
\addtolength{\textheight}{1.4in}
\addtolength{\oddsidemargin}{-.2in}
\addtolength{\textwidth}{.4in}
\pagestyle{empty}

\newcommand{\ds}{\displaystyle}
\newcommand{\ts}{\textstyle}

\newcommand{\Z}{\mathbb{Z}}
\newcommand{\R}{\mathbb{R}}
\newcommand{\C}{\mathbb{C}}
\newcommand{\EE}{\mathcal{E}}
\newcommand{\tZ}{\tilde{\Z}}

\newcommand{\Sym}{\mathrm{Sym}}

\renewcommand{\circle}[1]{{\bigcirc\hspace{-.75em}#1}}

\begin{document}

\footnote[0]{Notes by Peter Magyar \ {\tt magyar@math.msu.edu}}
\vspace{-1em}

\centerline{\bf Math 133\hspace{1.5em} \hfill{\Large\bf
Series}\hfill Stewart \S11.2}
\vspace{1em}

\noindent 
{\bf Sequences and series.} To any sequence $\{a_n\}_{n=1}^\infty$ 
we associate another sequence $\{s_n\}_{n=1}^\infty$, 
called the {\it series} of sums of $\{a_n\}$, defined by:
$$
s_n \ =\ a_1 + a_2 + \cdots + a_n
\ =\ \sum_{i=1}^n a_i\,.
$$
That is, $s_1=a_1$, \ $s_2=a_1+a_2$,  \ $s_3=a_1+a_2+a_3$, \ and in general $s_n$ is the sum of the first $n$ entries of $\{a_n\}$. The sigma notation $\sum_{i=1}^n a_i$ is a 
convenient shorthand for taking each 
integer value $i=1, 2,\ldots, n$,
substituting it into the expression $a_i$, and adding all the resulting quantities (see \S4.1 Pt 2). 
We can change the index letters arbitrarily: $s_N=\sum_{n=1}^N a_n$ defines the same series as $s_n=\sum_{i=1}^n a_i$.
\\[.5em]
{\sc example:} 
If $a_n= \frac1n$, then $s_1=1$, \ $s_2=1+\frac12=\frac32$, \ $s_3=1+\frac12+\frac13=\frac{11}6$, 
etc. The general term $s_n=\sum_{i=1}^n \frac1i = \frac11+\frac12+\frac13+\cdots+\frac1n$
has no elementary formula, which is typical for series, even when $a_n$ is quite simple.
Geometrically, $s_n$ is the area under the
bar graph of $\{a_n\}$ above the interval $[0,n]$.
\\[1em]
\centerline{
\includegraphics[width=2.3in]
{11-2-Bar-Graph.png} 
}
\vspace{-.5em}

\noindent
We can also take the infinite sum, which is defined as a limit of finite sums:
$$
\sum_{i=1}^\infty a_i
\ =\ \lim_{n\to\infty} \sum_{i=1}^n a_i
\ =\ \lim_{n\to\infty} s_n\,.
$$
This limit, converging or diverging, is the total area under the bar graph of $\{a_n\}_{n=1}^\infty$.
\\[.5em]
{\sc example:} The general purpose of a series is to express a complicated quantity as an infinite sum of simple quantities, $s=\sum_{i=1}^\infty a_i$, so that the finite sums $s_n =\sum_{i=1}^n a_i$ are approximations. A familiar example of this is decimal notation: an irrational real number is a complicated quantity with an infinite amount of detail in its digits, which are equivalent to the sum of a certain series. 

Given a sequence of digits $\{d_n\}_{n=1}^\infty$ with $d_n\in\{0,1,\ldots, 9\}$,
we have the number:
$$
s \ =\ 0.d_1d_2d_3\cdots\ =\  \frac{d_1}{10^1} + \frac{d_2}{10^2} + \frac{d_3}{10^3} + \cdots
\ =\ \sum\limits_{n=1}^\infty \frac{d_n}{10^n}\,. 
$$
By definition, an infinite decimal is the limit of its finite decimal approximations,
the number approached as we add more digits.
A trivial example is the repeating decimal $s=0.999\cdots$, which
clearly gets as close as desired to 1 as we take more
digits, so $s=1$. For a more complicated pattern of digits, the series converges to some complicated real number (see \S11.4).
\\[.5em]
{\sc example:} We will eventually (\S11.9, 11.10) develop powerful methods to write familiar numbers and functions as infinite series. Two outstanding formulas are:
$$
\pi \ =\ 4(1 - \tfrac13 + \tfrac 15 -\tfrac17+\cdots)\,,
\qquad\ \ 
\sin(x) \ =\
x -\tfrac1{3!}x^3 +\tfrac1{5!}x^5 -\tfrac1{7!}x^7 +\cdots.
$$
It is formulas like these that allow machines to compute complicated transcendental quantities using only the four arithmetic operations (which are all that can easily be  built into a logic circuit).
\\[1em]
{\bf Geometric progression on a chessboard.} 
A classic puzzle says: if we put one kernel of wheat on the first square of a chessboard, then two kernels on the second square, then four on the third square, and we keep doubling until the $64$th  square, how many kernels on the whole board?

We start with the sequence $a_n$ = number of kernels on the $n$th square, defined recursively by 
$a_1=1$, \ $a_n=2a_{n-1}$, which leads to the explicit
formula $a_n=2^{n-1}$.
This is called an exponential sequence or {\it geometric sequence}.\footnote{Geometric progression: 2 = length of a segment, $2^2 =$ area of a square, $2^3 =$ volume of cube.}
The associated {\it geometric series} is $s_n$ = total number of kernels on the first $n$ squares:
$$
s_n \ =\  \sum_{i=1}^n 2^{i-1}
\ =\ 
1 + 2^1 + 2^2 + \cdots + 2^{n-2} + 2^{n-1}.
$$
Surprisingly, we can find a simple formula for $s_n$ as follows:
\\[-.4em]
$$\begin{array}{rcccccccccc}
2 s_n &=&& &2^1&+& 2^2&+\cdots + &2^{n-1}&+&2^n\\[.5em]
-s_n&=& -1 & - &2^1 & - &2^2 &-\cdots- &2^{n-1}&
\end{array}$$
Adding these, the two sides become:
$$(2-1)s_n \ =\ 2a_n-a_1 \ =\ 2^n-1
\qquad\Longrightarrow\qquad
s_n \ =\ \frac{2^n-1}{2-1} \ =\ 2^n-1\,.
$$
Therefore, the answer to our puzzle is $s_{64}=2^{64}-1$ kernels, which
is enough wheat to fill the bowl of Spartan Stadium.
\\[.5em]
Now we change the problem: we put one ounce of gold on the first square, half an ounce on the second square,
a quarter ounce on the third square,
and so on until the 64th.
What is the total weight of gold on the board?

The weight on the $n$th square is given by another geometric
(i.e.~exponential) sequence $a_n=\left(\frac12\right)^{\!n-1}$, and the total weight on the first $n$ squares is the geometric series
$s_n=\sum_{i=1}^n \left(\frac12\right)^{\!i-1}$.
We use the same trick as before:
$$\begin{array}{rcccccccccc}
s_n &=&1& + &\frac12&+& \left(\frac12\right)^{\!2}&+\cdots + &\left(\frac12\right)^{\!n-1}\\[.5em]
-\frac12s_n&=& & - &\frac12 & - &\left(\frac12\right)^{\!2} &-\cdots- &\left(\frac12\right)^{\!n-1}&-& \left(\frac12\right)^{\!n}
\end{array}$$
\\[-.8em]
$$
(1-\tfrac12)s_n  \ =\ a_1 -\tfrac12 a_n\ =\ 1-\left(\tfrac12\right)^{\!n}
\qquad\Longrightarrow\qquad
s_n \ =\ \frac{\ \ 1-\left(\frac12\right)^{\!n}}{1-\frac12} \,.
$$
Thus, the total weight is $s_{64} = \frac{\ \ 1-\left(\frac12\right)^{\!64}}{1-\frac12}$. Since
$\left(\frac12\right)^{\!64}$ is a tiny, negligeable quantity,
this is very close to \smash{$\frac{1}{1-\frac12}=2$,} meaning the 
first square has just about the same total
weight as the other 63 squares.
In fact, adding more squares would barely change 
the total, since the limit is:
$$
\sum_{i=1}^\infty \left(\tfrac12\right)^{\!i-1} 
\ =\ 
\lim_{n\to\infty} s_n
\ =\
\lim_{n\to\infty} \frac{1-\left(\frac12\right)^{\!n}}{1-\frac12}
\ = \ \frac{1-0}{1-\frac12}\ =\ 2\,.
$$
\\[1em]
{\bf Geometric sequences and series.} 
A general geometric sequence starts with an initial value $a_1=c$, and subsequent terms are multiplied by the ratio $r$,
so that $a_n = ra_{n-1}$; explicitly, $a_n = cr^{n-1}$.
The same trick as above gives a formula for the corresponding
geometric series. We have $s_n-rs_n = c-cr^n$, so:
$$
\boxed{
\ \ 
s_n \ =\
\sum_{i=1}^n cr^{i-1}
\ =\ 
c + cr + cr^2 +\cdots +cr^{n-1}
\ =\ 
c\,\frac{\ 1-r^n}{1-r}\,. \ \
}$$
\\[-.2em]
(Notice that the power $r^n$ is one larger than in the last term $cr^{n-1}$.)
This ingenious formula is known as the {\it sum of a finite geometric series};
the limit is the {\it sum of an infinite geometric series}:
$$
\boxed{\ \ 
\lim_{n\to\infty} s_n
\ =\
\sum_{i=1}^\infty c r^{i-1} 
\ = \ c\,\frac{1}{1-r}\,, \quad\text{provided $|r|<1$.}
\ \ }
$$
Of couse, the infinite series diverges if $|r|\geq 1$.
These formulas are needed again and again in practical problems, especially
those involving financal interest rates.
\\[1em]
{\bf Manipulating sigma notation.} The geometric series formulas allow us to evaluate any series whose terms involve
only exponential functions like $2^n$ or $2^{-n}$, but not power functions like $n$ or $n^2$.
To evaluate, we rearrange and manipulate the terms into
the form of the geometric sequences which we know.
\\[.5em]
{\sc example:} Evaluate the finite sum \
$
\ds \sum_{i=1}^n \frac{2^i-1}{3^{i+1}}
$,
\ and the infinite sum\ \
$
\ds\sum_{i=1}^\infty \frac{2^i-1}{3^{i+1}}\,.
$
First, we work this out in dot-dot-dot notation:
$$
\frac{2^1-1}{3^{1+1}} + \cdots + \frac{2^n-1}{3^{n+1}} 
\ =\
\left(\frac{2^1}{3^{1+1}} + \cdots + \frac{2^n}{3^{n+1}}\right) 
- 
\left(\frac{1}{3^{1+1}} + \cdots + \frac{1}{3^{n+1}}\right)
$$
$$
\ =\
\left(\frac{2}{\ 3^2}\frac{2^{0}}{3^{0}} + \cdots + \frac{2}{\ 3^2}\frac{2^{n-1}}{3^{n-1}}\right) 
- 
\left(\frac{1}{\ 3^2}\frac{1}{3^{0}} + \cdots + \frac{1}{\ 3^2}\frac{1}{3^{n-1}}\right)
$$
$$
\ =\
\left(\frac{2}{\ 3^2}\left(\frac23\right)^{\!0} + \cdots + \frac{2}{\ 3^2}\left(\frac23\right)^{\!n-1}\right) 
- 
\left(\frac{1}{\ 3^2}\left(\frac13\right)^{\!0} + \cdots + \frac{1}{\ 3^2}\left(\frac13\right)^{\!n-1}\right) 
$$
$$
\ =\
\tfrac2{\ 3^2}\,\frac{\ \ 1-\left(\frac23\right)^{\!n}}{1-\frac23} 
\ \ -\ \ 
\tfrac1{\ 3^2}\,\frac{\ \ 1-\left(\frac13\right)^{\!n}}{1-\frac13}
$$
Here we factored out $\frac2{\ 3^2}$ and $\frac{1}{3^2}$ so the remaining factor would be $r^{i-1}$ for some $r$ and $i=1,\ldots,n$, which we can then evaluate using the geometric series formula.

This computation can be written more compactly in sigma notation:

$$\begin{array}{rcl}
\sum\limits_{i=1}^n \frac{2^i-1}{3^{i+1}}
&=&
\sum\limits_{i=1}^n \left(\frac{2^i}{3^{i+1}}-\frac{1}{3^{i+1}}\right)
\\[1.2em]
&=&
\sum\limits_{i=1}^n \frac{2^i}{3^{i+1}}\ -\ \sum\limits_{i=1}^n\frac{1}{3^{i+1}}
\\[1.2em]
&=&
\sum\limits_{i=1}^n \frac{2}{\ 3^2}\frac{2^{i-1}}{3^{i-1}}\ -\ 
\sum\limits_{i=1}^n\frac1{\ 3^2}\frac{1}{3^{i-1}}
\\[1.2em]
&=&
\sum\limits_{i=1}^n \frac{2}{\ 3^2}\left(\frac23\right)^{\!i-1}\ -\ 
\sum\limits_{i=1}^n\frac{1}{\ 3^2}\left(\frac13\right)^{\!i-1}
\\[1.2em]
&=&
\frac2{\ 3^2}\,\dfrac{\ \ 1-\left(\frac23\right)^{\!n}}{1-\frac23} 
\ \ -\ \ 
\frac1{\ 3^2}\,\dfrac{\ \ 1-\left(\frac13\right)^{\!n}}{1-\frac13}\,.
\end{array}$$
To get the infinite sum, we just remove the terms $(\frac23)^n$ and $(\frac13)^n$,
since these go to zero as $n\to\infty$.
\\[1em]
{\bf Repeating decimals.} An important application of geometric series is to 
write repeating infinite decimals as fractions. For example, consider:
$$
s \ =\ 0.0626262\cdots  \ =\  0.0\overline{62}
$$
We can apply the ratio-multiplication trick to cancel the infinite tail of digits:
$$
s\ =\  0.0626262\cdots, \qquad
\tfrac{1}{100}s \ =\ 0.0006262\cdots 
$$
$$
(1-\tfrac{1}{100})s \ = \ 0.062 
\qquad\Longrightarrow\quad
\tfrac{99}{100}s  \ =\ \tfrac{62}{1000}
\qquad\Longrightarrow\quad
s  \ =\ \tfrac{62}{990}\,.
$$

It is no coincidence that we can use the same trick as before: in fact, this infinite decimal is the sum of a geometric series:
$$\begin{array}{rcl}
s &=& \frac{6}{10^2} + \frac{2}{10^3} + \frac{6}{10^4} + \frac{2}{10^5} 
 + \frac{6}{10^6} + \frac{2}{10^7} +\cdots
\ =\  \frac{62}{10^3} \ +\ \frac{62}{10^5} \ +\  \frac{62}{10^7}\ +\ \cdots
\\[1em]
&=&\ds
\sum\limits_{i=1}^\infty \frac{62}{10^{2i+1}}
\ =\
\sum\limits_{i=1}^\infty \frac{62}{10^1}\frac{1}{(10^2)^i}
\ =\
\sum\limits_{i=1}^\infty \frac{62}{10^3}\left(\!\frac1{100}\!\right)^{\!\!i-1}
\\[1.8em]
&=&\ds
\frac{62}{10^3}\,\dfrac{1}{1-\frac{1}{100}}
\ \ =\ \
\frac{62}{1000}\,\frac{100}{99}
\ \ =\ \
\frac{62}{990}.
\end{array}$$

\noindent
Another example:
$
1.5626262\cdots \ =\ 1.5 + 0.0626262\cdots 
\ =\ \frac{15}{10} +\frac{62}{990} \ =\ \frac{1547}{990}.
$
\\[.6em]
{\small
We can carry out such reasoning for any infinite decimal which starts with arbitrary digits,
then becomes repeating. Conversely, for any fraction $a/b$, the long division 
$b\,)\hspace{-.35em}\overline{\ \, \vphantom{A}a\,}$
will produce a repeated remainder after at most $b{-}1$ steps, 
leading to a repeating decimal with a period of at most $b{-}1$.
E.g.~$1/7 = 0.\overline{142857}$ has remainders $3,2,6,4,5,1,3,2, \ldots.$

Thus, any infinite decimal represents a real number, and the repeating decimals represent 
precisely the rational numbers (fractions)!
For example, since we know $\sqrt 2 = 1.4142135623730950488\cdots$ is an irrational number, not equal to any fraction, its decimal digits will never settle into an infinitely repeating pattern.
}
\\[1em]
{\bf Geometric power series.} We can take the ratio in a geometric series to be a variable, obtaining a function called a {\it power series}:
$$
g(x)\ =\ \sum_{n=0}^\infty x^n \ =\ 1 + x+ x^2 + x^3+\cdots
$$
(Traditionally, the index starts at $n=0$, so the first term is $x^0=1$.)
By definition, this means that for the input $x=r$, the output is:
$$
g(r) 
\ =\
\sum_{n=0}^\infty r^n
\ =\ 
\lim_{N\to\infty} \sum_{n=0}^N r^n\,,
$$
whenever this limit exists.
Our formula for the sum of a geometric series can be 
rewritten: {$\sum\limits_{n=0}^\infty cr^n = c\,\frac{1}{1-r}$} for $|r|<1$.
This means:
\\[-.4em]
$$\boxed{\quad
g(x) \ =\ \sum_{n=0}^\infty x^n
\ =\ \left\{\begin{array}{cl}
\dfrac1{1-x} & \text{if $|x|<1$}
\\[1em]
\text{undefined} &\text{if $|x|\geq1$}.
\end{array}\right.
\quad}$$
\\[-.2em]
The set of $x$ for which the series converges is called the {\it interval of convergence}:
in this case it is $-1<x<1$, i.e.~$x\in(-1,1)$.

Another example:
$$
g(x) 
\ =\
\sum_{n=0}^\infty 3^{n-1}x^{n+1}
\ =\
\sum_{n=0}^\infty \frac{x}{3} (3x)^n
\ =\
\frac{x}{3}\,\frac{1}{1-3x} 
\ = \
\frac{x}{3-9x}\,,
$$
provided $|3x|<1$. The interval of convergence is thus $x\in(-\frac13,\frac13)$.
\\[1em]
{\bf Testing convergence.} We usually cannot find a neat formula for the
sum of an infinite series $\sum_{i=1}^\infty a_i=\lim_{n\to\infty}\sum_{i=1}^n a_i$,
but we still wish to know whether the series converges to some finite value. 
The most obvious way to analyze this is:
\begin{quote}
{\it Nth Term Non-vanishing Test:} 
If $\lim\limits_{n\to\infty} a_n \neq 0$, then the series $\sum_{i=1}^\infty a_i$ diverges. 
If $\lim\limits_{n\to\infty} a_n = 0$, then the series might converge or diverge.
\end{quote}
%
{\sc examples:} Use the Non-Vanishing Test to detect divergence.

\begin{itemize}

\item  $\sum_{n=0}^\infty 3^{n-1}x^{n+1}$ from the previous example.
For the Non-vanishing Test, we want to know if the limit of the terms vanishes,
$\lim\limits_{n\to\infty} 3^{n-1}x^{n+1} \stackrel{?}= 0$.
Here $3^{n-1}x^{n+1} = (3x)^{n-1}x^2$, and $x$ is fixed while $n$ gets bigger.
We see the limit is non-zero provided $|3x|\geq1$, or $|x|\geq\frac13$: in this case the series diverges.

For $|x|<\frac13$, the terms do approach zero, but in this case
the Non-vanishing Test cannot determine convergence or divergence  
(though we know from our analysis of geometric series that it really does converge).
 
\item $\sum_{n=1}^\infty \frac1n = 1 + \frac12 + \frac 13 +\frac14+\cdots$. 
Here $\lim\limits_{n\to\infty} a_n = \lim\limits_{n\to\infty} \frac1n = 0$, 
so the Non-vanishing Test 
cannot determine whether the series converges.\footnote{
We will see in \S11.3 that it diverges by the Integral Test.
Or estimate by grouping terms:
$$
1 + \underbrace{\tfrac12}_1 + \underbrace{\tfrac13 + \tfrac14}_2 + \underbrace{\tfrac15 + \tfrac16+\tfrac17+\tfrac18}_4 +\cdots
> 1 + \tfrac12 + 2\cdot\tfrac14 + 4\cdot\tfrac18 +\cdots=1+ \tfrac12+ \tfrac12+ \tfrac12+\cdots. 
$$
\\[-3em]
}

\end{itemize}

\end{document}

\item  $\sum_{n=1}^\infty (-1)^n n^2 = 1^2 - 2^2 + 3^2 - 4^2 + \cdots$.
Clearly $\lim_{n\to\infty} n^2 =\infty\neq 0$, so the series diverges by the Non-vanishing Test.
This is intuitively clear: instead of zeroing in on a limiting value, each additional term makes the sum oscillate wildly up or down.

\item
$\sum_{n=1}^\infty (-1)^n  = 1 - 1 + 1 - 1 + \cdots$.
The limit of the sequence $\lim\limits_{n\to\infty} (-1)^n$ does not exist, 
so the series diverges. Indeed, the finite sum
$s_N=\sum_{n=1}^N (-1)^n = 1$ if $N$ is odd and $0$ if $N$ is even, 
and $\{s_N\}_{N=1}^\infty = 1,0,1,0,\ldots$ does not converge to a single value.


%%%%%%%%%%%  Picture  %%%%%%%%%%
\\[0em]
\centerline{
%\includegraphics[width=2in]
{1-8-Discont-iv.png}
\qquad \qquad 
%\includegraphics[width=2in]
{1-8-Discont-v.png} 
}
\vspace{0em}

\noindent
%%%%%%%%%%%%%%%%%%%%%%%%%

%%%%%%%%%%%  Table  %%%%%%%%%%
In fact, we get the following derivatives:
$$\begin{array}{|c||c|c|c|c|c|c|}
\hline &&&&&& \\[-.9em] 
f(x) & \sin(x) & \cos(x) & \tan(x)  & \sec(x) & \csc(x) & \cot(x)
\\[.2em] \hline &&&&&& \\[-.9em]
f'(x) &\cos(x) & -\sin(x) & \sec^2(x) & \tan(x)\sec(x) & -\cot(x)\csc(x) & -\csc^2(x)
\\[.1em] \hline
\end{array}$$
\\[1em]
%%%%%%%%%%%%%%%%%%%%%%%%%%



















